\documentclass[10pt]{article}
\usepackage[letterpaper,margin=0.74in]{geometry}
\usepackage{amsmath,amssymb,amsthm,mathtools}
\usepackage{graphicx}
\usepackage{microtype}
\usepackage[hidelinks]{hyperref}
\usepackage{xcolor}
\usepackage{booktabs}
\usepackage{enumitem}
\setlist{nosep,leftmargin=*}
\setlength{\parindent}{0pt}
\setlength{\parskip}{4pt}

\newtheorem{theorem}{Theorem}
\newtheorem{lemma}{Lemma}
\newcommand{\R}{\mathbb R}
\newcommand{\A}{\mathbf A}
\newcommand{\M}{\mathbf M}
\newcommand{\Hh}{\mathbf H}
\newcommand{\supp}{\operatorname{supp}}
\newcommand{\sigmaerr}[2]{\sigma_{#1\A,\M}(#2)}

\definecolor{resultblue}{RGB}{235,245,252}
\definecolor{resultborder}{RGB}{38,105,145}
\newcommand{\resultbox}[1]{%
  \begin{center}\fcolorbox{resultborder}{resultblue}{%
  \begin{minipage}{0.93\linewidth}\textbf{Result.} #1\end{minipage}}\end{center}}

\title{A five-dimensional counterexample to greedy suboptimality for sparse approximation triples}
\author{Solution packet for arXiv:1506.03549, Appendix B}
\date{13 August 2026}

\begin{document}
\maketitle
\vspace{-1.5em}

\resultbox{The greedy estimate (B.4) in arXiv:1506.03549 is false for
general sparse approximation triples, even in dimension five and even when
every signal is compressible.  All greedy choices in the example are unique.
For a family $x_\varepsilon$, the two-step greedy/best-error ratio is
$(1+\varepsilon)/\varepsilon\to\infty$; at three steps the best error is zero
but the greedy error remains positive.}

\section*{1. The question}

Sun and Tang define, for a sparse approximation triple
$(\A,\M,\Hh)$, the greedy iterates
\[
 x^{k+1}=x^k+\operatorname*{argmin}_{a\in\A}
 \|x-x^k-a\|_{\M},\qquad x^0=0.
\]
Their Appendix B asks whether there is a constant $C>0$ such that
\[
 \|x^k-x\|_{\M}\leq C\,\sigmaerr{k}{x},\qquad
 \sigmaerr{k}{x}:=\inf_{z\in k\A}\|x-z\|_{\M},                 \tag{1}
\]
even for compressible $x$ \cite{source}.  Here is the exact question from
page 27 of the official arXiv PDF.

\begin{center}
\includegraphics[width=0.93\linewidth]{source_question_crop.png}
\end{center}
\vspace{0.1em}
\begin{center}\small\emph{Official source evidence: arXiv:1506.03549, p.~27.}\end{center}

\section*{2. Proof intuition}

Coordinate subspaces turn this greedy algorithm into weighted greedy set
cover.  In $\ell_1$, the best approximation on a coordinate support simply
keeps the coordinates on that support.  Thus a greedy step chooses the
allowed support containing the largest remaining $\ell_1$ mass, whereas a
best $k$-term approximation chooses $k$ supports whose \emph{union} captures
the most mass.  Overlapping supports can lure greedy into covering two large
coordinates that the optimal pair would cover separately.

The only delicate issue is whether such a support family satisfies every
axiom imposed on a sparse approximation triple.  It does in finite
dimension: coordinate projections are simultaneously best in $\ell_1$ and
$\ell_2$, both required norm-splitting identities are exact, and all closed
sets are proximinal.

\section*{3. The fixed triple}

Let
\[
 \M=(\R^5,\|\cdot\|_1),\qquad
 \Hh=(\R^5,\|\cdot\|_2).
\]
For $S\subset\{1,\ldots,5\}$, write
$E_S=\{z\in\R^5:\supp z\subset S\}$.  Define
\[
 \begin{array}{c@{\qquad}c@{\qquad}c@{\qquad}c}
 S_1=\{1,2\},&S_2=\{3,4\},&S_3=\{1,3\},&S_4=\{5\},
 \end{array}
\]
and $\A=E_{S_1}\cup E_{S_2}\cup E_{S_3}\cup E_{S_4}$.

\begin{lemma}\label{lem:triple}
$(\A,\M,\Hh)$ is a sparse approximation triple in the sense of the source.
Moreover $3\A=\R^5$.
\end{lemma}

\begin{proof}
We verify the five source axioms.
\begin{enumerate}[label=(\roman*)]
\item The inclusions $\A\subset\M\subset\Hh$ are continuous:
$\|z\|_2\leq\|z\|_1$, and on every $E_{S_i}$,
$\|z\|_1\leq\sqrt2\|z\|_2$.
\item $\A$ is a finite union of closed subspaces.  Every nonempty closed
subset of finite-dimensional $\M$ is proximinal: a minimizing sequence for
distance from a fixed point is bounded and has a convergent subsequence.
\item If $P_S$ denotes coordinate projection, then $P_Sx$ is the unique
best approximation to $x$ from $E_S$ in both $\ell_1$ and $\ell_2$.
\item Coordinate separation gives exactly
\[
 \|x\|_1=\|P_Sx\|_1+\|x-P_Sx\|_1,
 \quad
 \|x\|_2^2=\|P_Sx\|_2^2+\|x-P_Sx\|_2^2.
\]
\item Since $S_1\cup S_2\cup S_4=\{1,\ldots,5\}$,
$E_{S_1}+E_{S_2}+E_{S_4}=\R^5$.  Hence $3\A=\Hh$, which is stronger than
the required density of $\bigcup_k k\A$.
\end{enumerate}
\end{proof}

For any residual $r$ and allowed support $S_i$,
\[
 \inf_{a\in E_{S_i}}\|r-a\|_1
 =\|r-P_{S_i}r\|_1
 =\|r\|_1-\sum_{j\in S_i}|r_j|.                        \tag{2}
\]
Thus the greedy support is exactly the one capturing the greatest remaining
coordinate mass.

\section*{4. The counterexample}

Fix $0<\varepsilon<1$ and put
\[
 x_\varepsilon=(2,\tfrac32,2,1,\varepsilon).
\]
The captured masses at the first three greedy steps are:
\[
\begin{array}{c@{\quad}c@{\quad}cccc@{\quad}c}
 k & \text{residual} & S_1&S_2&S_3&S_4&\text{unique choice}\\
\midrule
0&(2,\frac32,2,1,\varepsilon)&\frac72&3&4&\varepsilon&S_3\\
1&(0,\frac32,0,1,\varepsilon)&\frac32&1&0&\varepsilon&S_1\\
2&(0,0,0,1,\varepsilon)&0&1&0&\varepsilon&S_2
\end{array}
\]
The inequalities are strict because $0<\varepsilon<1$.  Hence the two- and
three-step greedy residuals are
\[
 r^2=(0,0,0,1,\varepsilon),\qquad
 r^3=\varepsilon e_5,                                   \tag{3}
\]
so their errors are $1+\varepsilon$ and $\varepsilon$.

\begin{theorem}\label{thm:main}
For the fixed sparse approximation triple of Lemma~\ref{lem:triple}, no
constant $C$ can satisfy (1), even when restricted to compressible signals
and even under unique greedy choices.
\end{theorem}

\begin{proof}
The pair $S_1,S_2$ covers coordinates $1,2,3,4$, so
$\sigmaerr{2}{x_\varepsilon}\leq\varepsilon$.  Every other union of two
allowed supports misses at least one of those four coordinates, whose weight
is at least one.  Therefore
\[
 \sigmaerr{2}{x_\varepsilon}=\varepsilon,qquad
 \frac{\|x^2-x_\varepsilon\|_1}{\sigmaerr{2}{x_\varepsilon}}
 =\frac{1+\varepsilon}{\varepsilon}\longrightarrow\infty.   \tag{4}
\]
This already rules out a uniform $C$ while the denominator is positive.

There is also a pointwise exact-sparsity failure.  Because
$x_\varepsilon\in E_{S_1}+E_{S_2}+E_{S_4}\subset3\A$,
\[
 \sigmaerr{3}{x_\varepsilon}=0,
 \qquad \|x^3-x_\varepsilon\|_1=\varepsilon>0              \tag{5}
\]
by (3).  Hence no finite constant works even for a fixed $x_\varepsilon$.

Finally, every $x\in\M$ is compressible in the source's sense.  Indeed,
$3\A=\M$ implies $\sigmaerr{k}{x}=0$ for $k\geq3$; for any $\alpha>0$, a
finite enlargement of the constant handles $k=1,2$ in
$\sigmaerr{k}{x}\leq C_xk^{-\alpha}$.  Thus (4)--(5) occur entirely inside
the requested compressible class.
\end{proof}

\section*{5. Structural meaning and reproducibility}

More generally, any finite family of coordinate supports whose finite unions
span $\R^d$ gives a sparse approximation triple with
$\M=\ell_1^d$, $\Hh=\ell_2^d$.  Formula (2) identifies the source algorithm
with weighted greedy maximum coverage, while best approximation from $k\A$
is optimal weighted coverage by $k$ allowed supports.  Consequently a
positive theorem needs an additional structural hypothesis on the support
family; the five axioms alone cannot imply (1).

The included exact-rational script
\texttt{code/verify\_counterexample.py} enumerates all support unions,
confirms the three strict greedy choices, and verifies (4)--(5) for several
rational $\varepsilon$.  Run:
\begin{verbatim}
python3 code/verify_counterexample.py
\end{verbatim}
The expected final line is
\texttt{coordinate-group greedy counterexample: PASS}.  The script is an
audit aid; the proof above is exact and self-contained.

Exact-id and exact-phrase searches of the run indexes found no duplicate.  A
bounded public search through 2026-08-13 found the arXiv and published source
but no later answer to this exact sparse-approximation-triple question.  The
claim here is the negative result above, not a characterization of support
families for which a positive Lebesgue bound holds.

\begin{thebibliography}{9}
\bibitem{source}
Q. Sun and W.-S. Tang,
\emph{Nonlinear frames and sparse reconstructions in Banach spaces},
arXiv:1506.03549; J. Fourier Anal. Appl. 23 (2017), 1118--1152,
especially (5.3) and Appendix B,
\url{https://arxiv.org/abs/1506.03549}.
\end{thebibliography}

\end{document}
